\section{实验总结}

本次实验围绕 NEMU 模拟器和 Cache 行为仿真展开，主要完成了从程序运行、访存 trace 获取，到 Cache 模拟器设计、正确性验证以及不同 Cache 参数分析的完整实验流程。

首先，本实验通过香山仿真环境中的 AM 编译矩阵乘法 workload，并使用 NEMU 模拟器运行生成的二进制程序。通过开启内存访问日志功能，实验获得了程序运行过程中的内存访问序列 trace.log。相比直接在真实硬件上观察程序访存行为，NEMU 提供了一个更加可控、可重复、便于分析的实验环境，使得程序执行过程中的访存地址和读写行为能够被记录下来，为后续 Cache 行为分析提供了基础数据。

其次，本实验自行实现了 Cache 模拟器。该模拟器以 NEMU 生成的 trace.log 为输入，通过软件方式模拟 Cache 的地址映射、Tag 匹配、命中判断和替换过程。实验中实现了直接映射、组相联映射和全相联映射等不同映射方式，并支持 FIFO 和 LRU 两种替换算法。通过 micro-benchmarking 方法，本实验构造了若干组简单访存序列，并将人工推导的 Hit/Miss 结果与模拟器输出结果进行对比，从而验证了模拟器核心逻辑的正确性。

在 Cache 参数分析方面，实验结果表明，Cache 容量对命中率具有明显影响。对于矩阵规模为 $N=128$ 的普通矩阵乘法程序，当 Cache 容量较小时，命中率长期维持在较低水平；当 Cache 容量达到单个矩阵大小附近时，命中率出现明显跃升。这说明程序工作集大小与 Cache 容量之间的匹配关系，是影响 Cache 性能的重要因素。

同时，实验也分析了 Cache 块大小、相联度和替换算法的影响。Cache 块大小主要影响程序对空间局部性的利用程度；相联度主要影响冲突缺失。理论上，提高相联度可以减少多个活跃数据块竞争同一 Cache 位置所造成的冲突 Miss；但本实验中 2 路组相联命中率最高，继续提高到 4 路、8 路或全相联并没有带来单调提升，说明相联度的效果还会受到组数减少、地址分布和替换行为的共同影响。替换算法则影响 Cache 在容量有限时保留哪些数据块。在替换算法方面，LRU 从原理上更符合时间局部性思想，但本实验数据中 FIFO 与 LRU 的命中率非常接近，且 FIFO 在各组 Cache 容量下略高于 LRU。这说明替换算法的实际效果依赖于具体访存序列和 Cache 容量，当主要 Miss 来源是容量不足或程序访问模式对替换策略不敏感时，LRU 不一定表现出明显优势。但当程序 Miss 主要由容量不足或顺序扫描访问造成时，单纯改进替换算法带来的提升也会受到限制。

此外，通过不同矩阵规模的对比可以看出，随着矩阵规模增大，程序的工作集大小迅速增加，Cache 容量压力随之上升。对于普通三重循环矩阵乘法而言，矩阵 $A$ 的按行访问具有较好的空间局部性，矩阵 $C$ 的累加具有一定时间局部性，而矩阵 $B[k][j]$ 的按列访问会造成较大的访存步长，难以充分利用 Cache 块内的空间局部性。因此，非优化矩阵乘法的 Cache 性能不仅受硬件参数影响，也明显受代码访存模式影响。

在选做实验中，本实验进一步使用香山性能计数器统计了不同矩阵规模下的 DCache 访存行为。实验结果表明，随着矩阵规模从 $N=32$ 增大到 $N=128$，香山统计得到的 Cache 命中率从接近 $100\%$ 下降到约 $26.90\%$，说明在真实处理器微结构下，数据规模增大同样会显著增加 Cache 压力。与自行编写的 Cache 模拟器相比，香山性能计数器结果包含乱序执行、Load 重放、MissQueue 合并、多级 Cache 等硬件机制影响，因此二者在绝对数值上存在差异是合理的，但在“矩阵规模增大导致 Cache 性能下降”这一趋势上保持一致。

综上，本次实验加深了对 Cache 工作机制的理解。Cache 命中率并不是由单一因素决定的，而是由 Cache 容量、块大小、相联度、替换算法以及程序访存模式共同决定。对于数据密集型程序，除了调整 Cache 参数外，从程序层面改善局部性，例如循环交换、分块矩阵乘法等方法，也能够有效降低 Cache Miss，提高程序运行效率。
\newpage